\documentclass{ximera}
\newcommand{\blank}{\underline{\hspace{2.5cm}}}
\begin{document}

    \title{Homework - Exam Prep}
    
    \begin{abstract} Solve the following problems \textcolor{red}{BY HAND}.
    \\
    Due Date: none
    \end{abstract}
    \maketitle

\section{Projections}

\begin{exercise}
    State the definition of a projector, and the definition of an orthogonal projector.
\end{exercise}

% \begin{exercise}
%     Determine if the following matrices are a projector.  If so, also determine whether or not they are an orthogonal projector.
%     \begin{enumerate}
%         \item $\begin{bmatrix} 1/2 & 1/2 & 0 \\ 1/2 & 1/2 & 0 \\ 0&0&0 \end{bmatrix}$
%         \item $\begin{bmatrix} 1 & 1 & 0 \\ 0 & 0 & 0 \\ 0&0&0 \end{bmatrix}$
%         \item $\begin{bmatrix} 1 & 0 & 0 \\ 0 & 0 & 0 \\ 0&0&1 \end{bmatrix}$
%         \item $\begin{bmatrix} 1 & 1 & 0 \\ 1 & 1 & 0 \\ 0&0&1 \end{bmatrix}$
%     \end{enumerate}
% \end{exercise}

\begin{exercise}
Determine whether each of the following matrices is a projector. If so, determine whether it is an orthogonal projector. For each projector, describe geometrically what the matrix does to vectors in $\mathbb{R}^2$.

\begin{enumerate}
    \item
    \[
    P_1 =
    \begin{bmatrix}
    1 & 0 \\
    0 & 0
    \end{bmatrix}
    \]

    \item
    \[
    P_2 =
    \begin{bmatrix}
    1 & 1 \\
    0 & 0
    \end{bmatrix}
    \]

    \item
    \[
    P_3 =
    \begin{bmatrix}
    \frac{1}{2} & \frac{1}{2} \\
    \frac{1}{2} & \frac{1}{2}
    \end{bmatrix}
    \]

    \item
    \[
    P_4 =
    \begin{bmatrix}
    1 & 0 \\
    1 & 0
    \end{bmatrix}
    \]
\end{enumerate}
\end{exercise}


\begin{exercise}
Consider the matrix
\[
P =
\begin{bmatrix}
1 & 2 \\
0 & 0
\end{bmatrix}.
\]

\begin{enumerate}
    \item Show that $P$ is a projector.
    \item Find the range and nullspace of $P$.
    \item Describe geometrically what happens when $P$ acts on an arbitrary vector $x \in \mathbb{R}^2$.
    \item Is this an orthogonal projection? Explain geometrically why or why not.
    \item Sketch the projection of the vectors
    \[
    x_1 =
    \begin{bmatrix}
    1 \\ 1
    \end{bmatrix},
    \qquad
    x_2 =
    \begin{bmatrix}
    1 \\ -1
    \end{bmatrix}.
    \]
\end{enumerate}
\end{exercise}





% \begin{exercise}
% Consider the following two matrices:
% \[
% P_1 =
% \begin{bmatrix}
% 1 & 0 \\
% 0 & 0
% \end{bmatrix},
% \qquad
% P_2 =
% \begin{bmatrix}
% 1 & 1 \\
% 0 & 0
% \end{bmatrix}.
% \]

% \begin{enumerate}
%     \item Show that both matrices are projectors.
%     \item Determine whether either is an orthogonal projector.
%     \item Find the range and nullspace of each matrix.
%     \item Explain why both matrices project onto the same subspace, but produce different results for a general vector.
%     \item Sketch the geometric action of each matrix on the vector
%     \[
%     x =
%     \begin{bmatrix}
%     1 \\ 2
%     \end{bmatrix}.
%     \]
% \end{enumerate}
% \end{exercise}

\section{Conceptual Questions - SVD + QR}

\begin{exercise}
Suppose we want to project vectors in $\mathbb{R}^3$ onto the plane
\[
W = \operatorname{span}\left\{
\begin{bmatrix}
1 \\ 1 \\ 0
\end{bmatrix},
\begin{bmatrix}
0 \\ 1 \\ 1
\end{bmatrix}
\right\}.
\]
\begin{enumerate}
    \item Let $\tilde{A}$ be the matrix whose columns are the vectors whose span is $W$.  Use Gram-Schmidt to orthogonalize the columns.
    \item Now, using the dot product, create a matrix that, when applied to any vector $x \in \mathbb{R}^3$, projects it onto the subspace $W$.
    \item Verify that your matrix is a projector.
    \item Verify that your matrix is symmetric.
    \item Find the projection of
    \[
    b =
    \begin{bmatrix}
    1 \\ 0 \\ 1
    \end{bmatrix}
    \]
    onto $W$.
    % \item Explain geometrically why the residual $b-Pb$ is perpendicular to the plane $W$.
\end{enumerate}
\end{exercise}

\begin{exercise}
    Suppose a matrix $A$ has the SVD
    \begin{align*}
        A = \begin{bmatrix}  1 &0 &0 \\ 0& 1 & 0 \\ 0 & 0 &1\end{bmatrix} \begin{bmatrix} 3 &0 \\ 0 & 2 \\ 0 & 0\end{bmatrix} \begin{bmatrix}
            1 & 0 \\ 0 & 1
        \end{bmatrix}.
    \end{align*}
    Describe geometrically how $A$ transforms the unit circle.  Is this a full or reduced SVD formulation?  If it is the full SVD, write down the corresponding reduced SVD.  If it is a reduced SVD, generate a full SVD.
\end{exercise}

\begin{exercise}
    Explain the difference between classical and modified Gram-Schmidt.
\end{exercise}

\begin{exercise}
    Explain if/when QR factorization is needed when finding an SVD.
\end{exercise}

\section{Householder}

    
\begin{exercise}
Let
\[
x=\begin{bmatrix}3\\4\end{bmatrix}.
\]
Construct a Householder reflector $H$ that maps $x$ to a multiple of the first standard basis vector $e_1$.

\begin{enumerate}
    \item Find a vector $v$ such that
    \[
    H=I-2\frac{vv^T}{v^Tv}
    \]
    maps $x$ to a multiple of $e_1$.
    \item Construct the matrix $H$ explicitly.
    \item Verify that $Hx=\alpha e_1$ for some scalar $\alpha$.
    \item Verify that $H$ is symmetric and orthogonal.
    \item Explain geometrically what the Householder reflector does.
\end{enumerate}
\end{exercise}

\begin{exercise}
    Explain why a Householder reflector preserves the length of vectors.
\end{exercise}
\begin{exercise}
Consider the vector
\[
x=\begin{bmatrix}2\\-1\\2\end{bmatrix}.
\]

Construct a Householder reflector $H$ that maps $x$ to a multiple of $e_1$.

\begin{enumerate}
    \item Determine an appropriate vector $v$.
    \item Construct the Householder matrix $H$.
    \item Compute $Hx$.
    \item Verify that $H^TH=I$ and $H^2=I$.
    \item Explain why the Householder transformation preserves the length of vectors.
\end{enumerate}
\end{exercise}


% \begin{exercise}
% Consider the matrix
% \[
% A=
% \begin{bmatrix}
% 1&2\\
% 2&1\\
% 2&0
% \end{bmatrix}.
% \]

% Use a Householder reflector to introduce zeros below the first entry of the first column.

% \begin{enumerate}
%     \item Construct a Householder reflector $H_1$ such that
%     \[
%     H_1
%     \begin{bmatrix}
%     1\\2\\2
%     \end{bmatrix}
%     =
%     \begin{bmatrix}
%     \alpha\\0\\0
%     \end{bmatrix}.
%     \]
%     \item Compute $H_1A$.
%     \item Describe the structure of the resulting matrix.
%     \item Explain why applying an orthogonal transformation does not change the length of the columns of $A$.
% \end{enumerate}
% \end{exercise}


% \begin{exercise}
% Consider the matrix
% \[
% A=
% \begin{bmatrix}
% 1&2\\
% 2&1\\
% 2&0
% \end{bmatrix}.
% \]

% Use Householder reflectors to compute a QR factorization of $A$.

% \begin{enumerate}
%     \item Find a Householder reflector $H_1$ that introduces zeros below the first entry of the first column.
%     \item Find a second Householder reflector $H_2$ that introduces a zero below the diagonal in the second column of $H_1A$.
%     \item Compute
%     \[
%     R=H_2H_1A.
%     \]
%     \item Determine the orthogonal matrix $Q$ such that
%     \[
%     A=QR.
%     \]
%     \item Verify that $Q^TQ=I$ and $QR=A$.
% \end{enumerate}
% \end{exercise}


% \begin{exercise}
% Let $H$ be a Householder reflector defined by
% \[
% H=I-2\frac{vv^T}{v^Tv},
% \]
% where $v\neq 0$.

% \begin{enumerate}
%     \item Show that $H$ is symmetric.
%     \item Show that $H$ is orthogonal.
%     \item Show that $H^2=I$.
%     \item Show that $H$ has eigenvalue $-1$ in the direction of $v$.
%     \item Show that every vector perpendicular to $v$ is an eigenvector with eigenvalue $1$.
%     \item Explain geometrically why $H$ represents a reflection across the hyperplane perpendicular to $v$.
% \end{enumerate}
% \end{exercise}

    
\section{Condition Numbers, LU Decomposition}

\begin{exercise}
Consider the matrix
\[
A=
\begin{bmatrix}
1&0\\
0&\varepsilon
\end{bmatrix},
\qquad 0<\varepsilon<1.
\]

\begin{enumerate}
    \item Compute the $2$-norm condition number $\kappa_2(A)$.
    \item What happens to the condition number as $\varepsilon\to 0$?
    \item Explain geometrically what happens to a unit circle when it is transformed by $A$.
    \item Describe why solving $Ax=b$ becomes increasingly sensitive to perturbations as $\varepsilon$ approaches zero.
\end{enumerate}
\end{exercise}


\begin{exercise}
    Suppose you have the LU decomposition
    \begin{align*}
        A &= LU = 
        \begin{bmatrix}
            1 & 0 & 0\\
            2 & 1 & 0\\
            -1 & 3 & 1
        \end{bmatrix}
        \begin{bmatrix}
            2 & 1 & 1\\
            0 & 3 & 2\\
            0 & 0 & 4
        \end{bmatrix}.
    \end{align*}
    
    Use forward and back substitution to solve
    \begin{align*}
        Ax=b,
        \qquad
        b=
        \begin{bmatrix}
            4\\
            11\\
            9
        \end{bmatrix}.
    \end{align*}
    
    \begin{enumerate}
        \item First, solve $Ly=b$ using forward substitution.
        \item Next, solve $Ux=y$ using back substitution.
        \item Verify your solution by computing $Ax$.
    \end{enumerate}
\end{exercise}


\begin{exercise}
    What is the difference between an elementary matrix and an elimination matrix when creating an LU decomposition of a matrix?
\end{exercise}


\begin{exercise}
Consider the matrix
\[
A=
\begin{bmatrix}
\varepsilon & 1\\
1 & 1
\end{bmatrix},
\qquad 0<\varepsilon\ll 1.
\]

\begin{enumerate}
    \item Compute the $LU$ factorization of $A$ without pivoting, where $L$ is unit lower triangular.

    \item Examine the entries of $L$ and $U$ as $\varepsilon\to 0$. What happens to the size of the entries in these matrices?

    \item Compute the condition number of $L$ and the condition number of $U$.  Use this to explain why the factorization without pivoting may lead to numerical difficulties.

    \item Now perform partial pivoting and find matrices $P$, $L$, and $U$ such that
    \[
    PA=LU.
    \]

    \item Compare the entries of the two factorizations. In particular, compare the sizes of the entries in the two lower triangular matrices.

    \item Compute the condition number of $L$ and the condition number of $U$.  Use this to explain why partial pivoting produces a more numerically stable factorization in this example.

    \item Does partial pivoting change the condition number of $A$? Explain the difference between the conditioning of a matrix and the numerical stability of an algorithm.
\end{enumerate}
\end{exercise}


\begin{exercise}
     Suppose a matrix $A$ has the SVD
    \begin{align*}
        A = \begin{bmatrix}  1 &0 &0 \\ 0& 1 & 0 \\ 0 & 0 &1\end{bmatrix} \begin{bmatrix} 3 &0 \\ 0 & 2 \\ 0 & 0\end{bmatrix} \begin{bmatrix}
            1 & 0 \\ 0 & 1
        \end{bmatrix}.
        \end{align*}
        Compute its condition number.  Will it be ``hard" or ``easy" to solve a least squares problem $\min_x \|b-Ax\|_2$?
\end{exercise}

\section{Pseudocode}

\begin{exercise}

Let $A \in \mathbb{R}^{m \times n}$, where $m \geq n$, have full column rank.
The goal is to compute the QR decomposition $A = QR$ using
Householder reflectors, where $Q \in \mathbb{R}^{m\times m}$ is orthogonal
and $R \in \mathbb{R}^{m \times n}$ is upper triangular.

\begin{enumerate}
    \item Fill in each blank to complete the Householder QR algorithm.
Assume \texttt{A[k:m,k:n]} denotes the submatrix of rows $k,\dots,m$ and
columns $k,\dots,n$, \texttt{sign(x)} returns $+1$ if $x \ge 0$ and $-1$
otherwise, and \texttt{norm(x,2)} computes the Euclidean norm.

\begin{lstlisting}
Algorithm HouseholderQR(A):
    Input: A in R^(m x n), with m >= n
    Output: Q in R^(m x m), R in R^(m x n)

    for k = 1, 2, ..., n:

        x <- A[k:m, k]
        v <- x
        v[1] <- v[1] + __________ (1) __________
        v <- __________ (2) __________
        
        A[k:m, k:n] <- A[k:m, k:n] - __________ (3) __________
        Store v as v_k

    R <- A

    Q <- identity matrix of size m x m
    for k = __________ (4) __________:
        Q[k:m, :] <- Q[k:m, :] - __________ (5) __________

    return Q, R
\end{lstlisting}

\item Answer the following questions.
\begin{enumerate}[label=\arabic*.]
    \item Why does the algorithm choose the sign in blank (1) to match
    the sign of \texttt{x[1]}?
    \item The reflector $H_k = I - 2v_kv_k^T$ is never formed explicitly.
    Why not, and what is done instead?
    \item What is the effect of step $k$ on the entries of column $k$
    below the diagonal?
    \item Why is the loop in blank (4) run in that particular order?
\end{enumerate}

\item Give the leading-order flop count of the factorization phase
(not forming $Q$) in terms of $m$ and $n$.

% \item Name one advantage of Householder QR over Modified Gram--Schmidt in floating-point arithmetic.

\end{enumerate}

% \section*{Instructor Answer Key}
% \begin{enumerate}[label=\arabic*.]
%     \item $\operatorname{sign}(x_1)\,\|x\|_2$ \quad i.e.\ \texttt{sign(x[1]) * norm(x,2)}
%     \item $v/\|v\|_2$
%     \item $2\,v\,(v^T A[k{:}m,k{:}n])$
%     \item $n, n-1, \dots, 1$ (backward)
%     \item $2\,v_k\,(v_k^T Q[k{:}m,:])$
% \end{enumerate}
%
% (b) 1. Avoids subtractive cancellation in $x_1 - \|x\|_2$ (which can
% lose all accuracy when $x$ is nearly parallel to $e_1$).
% 2. Forming $H_k$ costs $O(m^2)$ storage and $O(m^2 n)$ work per
% application; applying it as $A - 2v(v^TA)$ is a rank-one update costing
% $O(mn)$.
% 3. They become zero (the reflector maps $x$ to $-\operatorname{sign}(x_1)\|x\|_2 e_1$).
% 4. $Q = H_1H_2\cdots H_n$, so applying $H_n$ first to $I$ and working
% outward to $H_1$ yields $Q$ with the stored vectors.
%
% (c) $2mn^2 - \tfrac{2}{3}n^3$.
%
% (d) Householder produces $Q$ that is orthogonal to machine precision
% regardless of the conditioning of $A$ (backward stable); MGS loses
% orthogonality roughly in proportion to $\kappa(A)$.

\end{exercise}

\begin{exercise}

Assume $A \in \mathbb{R}^{m\times n}$ with $m \gg n$ has full column rank.
The following algorithm computes the SVD by first reducing $A$ to a small
square matrix.

\begin{enumerate}
    \item Fill in the blanks using the word bank. Some words may be used
    more than once, or not at all.

\begin{center}
\fbox{
\begin{minipage}{0.85\textwidth}
\centering
QR \quad
upper triangular \quad
orthogonal \quad
transpose \quad
diagonal \quad
descending \quad
$U_R$ \quad
inverse
\end{minipage}
}
\end{center}

\begin{lstlisting}
Algorithm RSVD(A):

    # Step 1: Compute the reduced ________ factorization
    [Q, R] = qr(A)        # Q is m x n, R is n x n

    # Step 2: Compute the SVD of the small matrix
    #         R is n x n and ________
    [U_R, Sigma, V] = svd(R)

    # Step 3: Recover the left singular vectors
    U = Q * ________      # Q has orthonormal columns and
                          # this factor is n x n orthogonal

    # Step 4: Return the decomposition
    # Sigma is ________ with the singular values
    # in ________ order on its diagonal
    Return U, Sigma, V
\end{lstlisting}

\item Show that the output satisfies $A = U\Sigma V^T$, and explain why
$U$ has orthonormal columns.

\item Explain why this approach can be cheaper than computing the SVD of
$A$ directly when $m \gg n$.

\item Suppose instead that $A \in \mathbb{R}^{m\times n}$ is wide, with
$n \gg m$ and full row rank. In words, how would you modify the
algorithm?

\end{enumerate}

% \section*{Instructor Answer Key}
% (a) QR; upper triangular; $U_R$; diagonal; descending.
%
% (b) $A = QR = Q(U_R\Sigma V^T) = (QU_R)\Sigma V^T$. Since
% $U^TU = U_R^TQ^TQU_R = U_R^TU_R = I_n$, the columns of $U$ are orthonormal.
%
% (c) Computing the SVD of the full $m\times n$ matrix directly requires
% bidiagonalizing $A$ and then iterating while accumulating the $m\times n$
% matrix of left singular vectors, which has a large constant in the
% $O(mn^2)$ cost. In RSVD, the QR costs $O(mn^2)$ with a small constant,
% the SVD (including all iterations) is done on an $n\times n$ matrix
% ($O(n^3)$), and the only other $m$-dependent work is one $O(mn^2)$
% multiply $QU_R$. For $m \gg n$ the total is smaller.
%
% (d) Compute the QR factorization of $A^T = QR$ with $Q$ of size
% $n\times m$ and $R$ of size $m\times m$, so $A = R^TQ^T$ with $R^T$ lower
% triangular. Take the SVD $R^T = U\Sigma V_R^T$ and set $V = QV_R$.
% Then $A = U\Sigma V^T$.

\end{exercise}
\newpage

% \section*{Answer Key}

% \subsection*{Part 1: Pseudocode}

% The suggested answers for the blanks are:

% \begin{enumerate}
%     \item \texttt{transpose}
%     \item \texttt{descending}
%     \item \texttt{square root}
%     \item \texttt{eigenvectors}
%     \item \texttt{diagonal}
% \end{enumerate}

% The Gram matrix is $G = A^T A$. Its eigenvalues are
% nonnegative, and the singular values of $A$ are their
% nonnegative square roots.

% The right singular vectors are the normalized eigenvectors
% of $A^T A$.

% For a full SVD, $U$ and $V$ are orthogonal matrices, and
% $\Sigma$ is an $m \times n$ diagonal-shaped matrix. If
% $\Sigma$ is square and invertible, then
% $U = AV\Sigma^{-1}$ can be used. Zero singular values require
% special handling because $\Sigma^{-1}$ does not exist.

% \subsection*{Part 2: Multiple-Choice Answers}

% \begin{enumerate}[label=\arabic*.]
%     \item (b) The eigenvalues of $A^T A$ are the squared
%     singular values of $A$.
%     \item (a) The singular value is $\sqrt{9} = 3$.
%     \item (c) The columns of $V$ are the right singular vectors.
%     \item (a) A zero diagonal entry makes $\Sigma$ singular.
% \end{enumerate}

% \subsection*{Part 3: Challenge Solution}

% Given

% \[
% A =
% \begin{bmatrix}
% 3 & 0\\
% 0 & 2
% \end{bmatrix},
% \]

% we have

% \[
% A^T A =
% \begin{bmatrix}
% 9 & 0\\
% 0 & 4
% \end{bmatrix}.
% \]

% The eigenvalues of $A^T A$ are

% \[
% \lambda_1 = 9,
% \qquad
% \lambda_2 = 4.
% \]

% Therefore, the singular values are

% \[
% \sigma_1 = \sqrt{9} = 3,
% \qquad
% \sigma_2 = \sqrt{4} = 2.
% \]

% One valid decomposition is

% \[
% U =
% \begin{bmatrix}
% 1 & 0\\
% 0 & 1
% \end{bmatrix},
% \qquad
% \Sigma =
% \begin{bmatrix}
% 3 & 0\\
% 0 & 2
% \end{bmatrix},
% \qquad
% V =
% \begin{bmatrix}
% 1 & 0\\
% 0 & 1
% \end{bmatrix}.
% \]

% Indeed,

% \[
% U\Sigma V^T
% =
% \begin{bmatrix}
% 1 & 0\\
% 0 & 1
% \end{bmatrix}
% \begin{bmatrix}
% 3 & 0\\
% 0 & 2
% \end{bmatrix}
% \begin{bmatrix}
% 1 & 0\\
% 0 & 1
% \end{bmatrix}
% =
% \begin{bmatrix}
% 3 & 0\\
% 0 & 2
% \end{bmatrix}
% = A.
% \]

% \bigskip

% \noindent
% \textit{Instructor note:} This exercise illustrates the
% mathematical idea behind SVD using $A^T A$. In practical
% numerical computation, algorithms that avoid explicitly
% forming $A^T A$ are generally preferred because forming
% the Gram matrix can worsen numerical accuracy.


\end{document}